\documentclass[12pt]{article}
\usepackage[ukrainian,shorthands=off]{babel}   % shorthands=off: символ " не активний (потрібно для міток "..." у TikZ всередині \zadtask)
\usepackage{fontspec}
% --- НАЛАШТУВАННЯ ШРИФТІВ ---
% Шрифти Schoolbook-*.otf лежать у корені проєкту. Шлях підбирається автоматично:
% ../ (компіляція з папки теми), ./ (корінь проєкту), ../../ (підпапки теми 28); інакше -- TeX Gyre Schola.
\newcommand{\nmtSchoolbook}[1]{\setmainfont{Schoolbook-Regular.otf}[
    Path = #1,
    BoldFont = Schoolbook-Bold.otf,
    ItalicFont = Schoolbook-Italic.otf,
    BoldItalicFont = Schoolbook-BoldItalic.otf
]}
\IfFileExists{../Schoolbook-Regular.otf}{\nmtSchoolbook{../}}{%
 \IfFileExists{./Schoolbook-Regular.otf}{\nmtSchoolbook{./}}{%
  \IfFileExists{../../Schoolbook-Regular.otf}{\nmtSchoolbook{../../}}{%
   \setmainfont{texgyreschola}[Extension=.otf, UprightFont=*-regular, BoldFont=*-bold, ItalicFont=*-italic, BoldItalicFont=*-bolditalic]}}}
\usepackage[italic]{mathastext}
\usepackage[a4paper,margin=1.8cm,bottom=2cm,top=2cm]{geometry}
\usepackage{amsmath,amssymb}
\usepackage{tikz}
\usepackage{pgfplots}
\pgfplotsset{compat=1.18}
\usepgfplotslibrary{fillbetween}
\usetikzlibrary{calc,patterns,angles,quotes,intersections,babel,3d,shapes.symbols,shapes.geometric,decorations.pathreplacing,shadings,arrows.meta,hobby}
\usepackage{xcolor,array,fancyhdr,enumitem,multicol}
\usepackage{colortbl,multirow,diagbox,needspace}
\usepackage{tcolorbox}
\tcbuselibrary{skins,breakable}
\usepackage[hidelinks]{hyperref}
\usepackage[firstpage=false, text=@pvtr2525, scale=0.6, color=gray!12]{draftwatermark}

% --- АВТОСТИСКАННЯ: рисунок/сітка, ширша за колонку, масштабується до ширини колонки ---
\newsavebox{\nmtfitbox}
\newenvironment{nmtfit}{\begin{lrbox}{\nmtfitbox}}{\end{lrbox}%
  \ifdim\wd\nmtfitbox>\linewidth \resizebox{\linewidth}{!}{\usebox{\nmtfitbox}}\else\usebox{\nmtfitbox}\fi}
\let\nmtIncludegraphicsOrig\includegraphics
\renewcommand{\includegraphics}[2][]{\begin{nmtfit}\nmtIncludegraphicsOrig[#1]{#2}\end{nmtfit}}


\definecolor{mainGreen}{RGB}{34, 120, 64}
\definecolor{yearOrange}{RGB}{220, 100, 30}
\definecolor{yearcolor}{RGB}{220, 100, 30}
\definecolor{titleBg}{RGB}{225, 240, 220}
\definecolor{instrBg}{RGB}{255, 240, 220}
\definecolor{instrBorder}{RGB}{220, 150, 80}
\definecolor{headerblue}{RGB}{0, 102, 204}   % використовується в рисунках старої бази

\setlength{\parindent}{0pt}
\emergencystretch=1.5em

\pagestyle{fancy}
\fancyhf{}
\renewcommand{\headrulewidth}{0.4pt}
\renewcommand{\footrulewidth}{0.4pt}
\fancyhead[C]{\small\color{gray!80} База завдань НМТ 2022--2026 \textendash{} Тема 9}
\fancyhead[R]{\small\color{mainGreen}\textbf{\href{https://t.me/pvtr2525}{@pvtr2525}}}
\fancyfoot[L]{\small\color{gray!80}\href{https://t.me/pvtr2525}{ТГ: @pvtr2525}}
\fancyfoot[C]{\small\color{gray!80}\thepage}
\fancyfoot[R]{\small\color{gray!80}НМТ 2022--2026}

% --- ТАБЛИЦІ ВІДПОВІДЕЙ ---
% ширина клітинки = (ширина рядка - відступи) / 5: на всю сторінку це ~3 см, у minipage поруч із рисунком таблиця стискається
\newlength{\nmtcellw}
\newcommand{\nmtSetCell}{\setlength{\nmtcellw}{\dimexpr(\linewidth-12\tabcolsep-6\arrayrulewidth)/5\relax}}
\newcommand{\answerTable}[5]{%
\par\nopagebreak\vspace{0.15cm}
\noindent\nmtSetCell
\begin{tabular}{|*{5}{>{\centering\arraybackslash}m{\nmtcellw}|}}
\hline
\rule[-0.15cm]{0pt}{0.6cm}\textbf{А} & \rule[-0.15cm]{0pt}{0.6cm}\textbf{Б} & \rule[-0.15cm]{0pt}{0.6cm}\textbf{В} & \rule[-0.15cm]{0pt}{0.6cm}\textbf{Г} & \rule[-0.15cm]{0pt}{0.6cm}\textbf{Д} \\
\hline
\rule[-0.2cm]{0pt}{0.75cm}#1 & \rule[-0.2cm]{0pt}{0.75cm}#2 & \rule[-0.2cm]{0pt}{0.75cm}#3 & \rule[-0.2cm]{0pt}{0.75cm}#4 & \rule[-0.2cm]{0pt}{0.75cm}#5 \\
\hline
\end{tabular}
\par\vspace{0.25cm}
}
\newcommand{\answerTableTall}[5]{%
\par\nopagebreak\vspace{0.15cm}
\noindent\nmtSetCell
\begin{tabular}{|*{5}{>{\centering\arraybackslash}m{\nmtcellw}|}}
\hline
\rule[-0.2cm]{0pt}{0.7cm}\textbf{А} & \rule[-0.2cm]{0pt}{0.7cm}\textbf{Б} & \rule[-0.2cm]{0pt}{0.7cm}\textbf{В} & \rule[-0.2cm]{0pt}{0.7cm}\textbf{Г} & \rule[-0.2cm]{0pt}{0.7cm}\textbf{Д} \\
\hline
\rule[-0.45cm]{0pt}{1.2cm}#1 & \rule[-0.45cm]{0pt}{1.2cm}#2 & \rule[-0.45cm]{0pt}{1.2cm}#3 & \rule[-0.45cm]{0pt}{1.2cm}#4 & \rule[-0.45cm]{0pt}{1.2cm}#5 \\
\hline
\end{tabular}
\par\vspace{0.25cm}
}
\newcommand{\instructionBox}[1]{%
\par\vspace{0.3cm}
\begin{tcolorbox}[colback=instrBg, colframe=instrBorder, boxrule=0.6pt, arc=2pt,
    left=8pt, right=8pt, top=4pt, bottom=4pt]
\centering\small\bfseries #1
\end{tcolorbox}
\par\vspace{0.15cm}
}
\newcommand{\sectionTitle}[1]{%
\par\vspace{0.4cm}
\begin{tcolorbox}[colback=titleBg, colframe=titleBg, boxrule=0pt, arc=8pt,
    halign=center, fontupper=\Large\bfseries\color{mainGreen}]
\hyphenpenalty=10000 \exhyphenpenalty=10000 #1
\end{tcolorbox}
\par\vspace{0.3cm}
}
\newcommand{\task}[2]{\noindent\textbf{#1.}\ \ #2\par\vspace{0.2cm}}
% --- НМТ-СТИЛЬ ПОЛЕ ВІДПОВІДІ ---
\newcommand{\nmtAnswerBox}{%
\par\nopagebreak\vspace{0.2cm}\noindent
Відповідь:\ \
\begin{tabular}{@{}*{4}{|>{\centering\arraybackslash}p{0.8cm}}|@{\hspace{0.1cm}\textbf{,}\hspace{0.1cm}}*{3}{|>{\centering\arraybackslash}p{0.8cm}}|@{}}
\hline
\rule[-0.2cm]{0pt}{0.7cm} & & & & & & \\
\hline
\end{tabular}
\par\vspace{0.3cm}
}
% --- СІТКА ДЛЯ МАТЧИНГУ ---
\newsavebox{\nmtgridbox}
\newcommand{\matchingGrid}{%
    \sbox{\nmtgridbox}{\matchingGridRaw}%
    \ifdim\wd\nmtgridbox>\linewidth \resizebox{\linewidth}{!}{\usebox{\nmtgridbox}}\else\usebox{\nmtgridbox}\fi}
\newcommand{\matchingGridRaw}{%
    \begingroup
    \renewcommand{\arraystretch}{1.15}
    \setlength{\tabcolsep}{4pt}
    \footnotesize
    \begin{tabular}{r|c|c|c|c|c|}
         \multicolumn{1}{c}{} & \multicolumn{1}{c}{\textbf{А}} & \multicolumn{1}{c}{\textbf{Б}} & \multicolumn{1}{c}{\textbf{В}} & \multicolumn{1}{c}{\textbf{Г}} & \multicolumn{1}{c}{\textbf{Д}} \\ \cline{2-6}
         \textbf{1} & \phantom{X} & \phantom{X} & \phantom{X} & \phantom{X} & \phantom{X} \\ \cline{2-6}
         \textbf{2} & & & & & \\ \cline{2-6}
         \textbf{3} & & & & & \\ \cline{2-6}
    \end{tabular}
    \endgroup
}
% --- ДОДАТКОВІ МАКРОСИ ЗБІРНИКА ---
\newcommand{\taskBlock}[1]{%
\par\vspace{0.45cm}
\noindent{\color{mainGreen}\rule{\linewidth}{0.8pt}}\par\vspace{0.12cm}
\noindent{\small\bfseries\color{mainGreen}#1}\par\vspace{0.25cm}
}
\newcounter{zad}
\newcommand{\zadtask}[1]{\stepcounter{zad}\noindent\textbf{\thezad.}\ \ #1\par\nopagebreak\vspace{0.2cm}}
\newcommand{\zadnum}{\stepcounter{zad}\textbf{\thezad.}\ \ }
\newcounter{ans}
\newcommand{\ansitem}[1]{\stepcounter{ans}\noindent\textbf{\theans.}\ #1\par\vspace{0.07cm}}
% мітка року: нерозривна, притиснута праворуч; якщо не вміщається -- праворуч на новому рядку
\newcommand{\nmtyear}[1]{\unskip\nobreak\finalhyphendemerits=0 \hfill\penalty100\hskip1em\hbox{}\nobreak\hfill{\small\color{yearOrange}\mbox{\rule{0pt}{2.3ex}(НМТ~#1)}}}
% --- МАКРОС КОРОТКОГО РОЗВ'ЯЗКУ ---
\newcommand{\solution}[1]{%
\par\vspace{0.12cm}
\begin{tcolorbox}[colback=titleBg!45, colframe=mainGreen!55, boxrule=0.4pt, arc=2pt,
    left=6pt, right=6pt, top=3pt, bottom=3pt, breakable]
\footnotesize\textbf{Короткий розв'язок.}\ #1
\end{tcolorbox}
\par\vspace{0.12cm}
}
% Розв'язки й відповіді (є до завдань НМТ-2026): \showsolutionstrue -- показати, \showsolutionsfalse -- сховати
\newif\ifshowsolutions
\showsolutionsfalse
% --- ЗАГОЛОВКИ З ПУНКТАМИ КЛІКАБЕЛЬНОГО ЗМІСТУ ---
\setcounter{tocdepth}{2}
\newcommand{\chapterTitle}[1]{%
\clearpage\phantomsection\addcontentsline{toc}{section}{#1}%
\sectionTitle{#1}}
\newcommand{\typeTitle}[1]{%
\needspace{6\baselineskip}%
\phantomsection\addcontentsline{toc}{subsection}{\quad #1}%
\par\vspace{0.3cm}
\noindent{\large\bfseries\color{yearOrange}#1}\par\vspace{0.08cm}
\noindent{\color{yearOrange}\rule{\linewidth}{0.4pt}}\par\nopagebreak\vspace{0.2cm}}
\raggedbottom
% усередині одного завдання сторінка не розривається: нерозривний вертикальний проміжок
% і нерозривна межа абзаців (умова ніколи не відділяється від варіантів, рисунка чи сітки)
\newcommand{\nbvspace}[1]{\nopagebreak\vspace{#1}}
\newcommand{\nmtnobreak}{\ifvmode\penalty10000 \fi}
% --- ЄДИНИЙ МАКЕТ ЗАВДАНЬ НА ВІДПОВІДНІСТЬ ---
% мітка в колонці сталої ширини, текст -- у parbox: продовження довгого рядка
% вирівнюється під текстом, а не під міткою; відступ між пунктами однаковий скрізь
\newlength{\nmtMatchLab}\setlength{\nmtMatchLab}{1.6em}
\newlength{\nmtMatchSep}\setlength{\nmtMatchSep}{0.28cm}
\newcommand{\matchHead}[1]{\par\noindent\textit{#1}\par\nopagebreak\vspace{0.2cm}}
\newcommand{\matchItem}[2]{\par\noindent\makebox[\nmtMatchLab][l]{\textbf{#1}}%
\parbox[t]{\dimexpr\linewidth-\nmtMatchLab\relax}{\raggedright #2}\par\nopagebreak\vspace{\nmtMatchSep}}
\newcommand{\ansTheme}[1]{\par\vspace{0.25cm}\noindent{\bfseries\color{mainGreen}#1}\par\vspace{0.12cm}}
\newcommand{\ansType}[1]{\par\vspace{0.1cm}\noindent{\itshape\color{yearOrange}#1}\par\vspace{0.08cm}}

\begin{document}
\begingroup\setcounter{zad}{0}
% --- сумісність зі старою базою (діє, якщо тема не має власного визначення) ---
\providecommand{\trigStyle}[1]{#1}
\providecommand{\answerTableSmall}[5]{%
\begingroup\setlength{\tabcolsep}{2pt}\small\nmtSetCell
\begin{tabular}{|*{5}{>{\centering\arraybackslash}m{\nmtcellw}|}}
\hline
\rule[-0.2cm]{0pt}{0.6cm}\textbf{А} & \textbf{Б} & \textbf{В} & \textbf{Г} & \textbf{Д} \\
\hline
\rule[-0.4cm]{0pt}{0.9cm}#1 & \rule[-0.4cm]{0pt}{0.9cm}#2 & \rule[-0.4cm]{0pt}{0.9cm}#3 & \rule[-0.4cm]{0pt}{0.9cm}#4 & \rule[-0.4cm]{0pt}{0.9cm}#5 \\
\hline
\end{tabular}\endgroup}
\providecommand{\matchingLayout}[3]{%
    \noindent
    \begin{minipage}[t]{0.36\textwidth}\vspace{0pt}\raggedright
        #1
    \end{minipage}%
    \hfill
    \begin{minipage}[t]{0.43\textwidth}\vspace{0pt}\raggedright
        #2
    \end{minipage}%
    \hfill
    \begin{minipage}[t]{0.19\textwidth}
        \vspace{0pt}
        \begin{flushright}
        #3
        \end{flushright}
    \end{minipage}}
\providecommand{\matchingLayoutBottom}[2]{%
    \noindent
    \begin{minipage}[t]{0.48\textwidth}\vspace{0pt}\raggedright
        #1
    \end{minipage}%
    \hfill
    \begin{minipage}[t]{0.48\textwidth}\vspace{0pt}\raggedright
        #2
    \end{minipage}
    \vspace{0.5cm}
    \noindent
    \matchingGrid}
\providecommand{\answerListVertical}[5]{%
    \par\nopagebreak\vspace{0.2cm}
    \begin{itemize}[itemsep=0.4cm, leftmargin=1.5cm, labelsep=0.5cm]
        \item[\textbf{А}] #1
        \item[\textbf{Б}] #2
        \item[\textbf{В}] #3
        \item[\textbf{Г}] #4
        \item[\textbf{Д}] #5
    \end{itemize}
    \vspace{0.2cm}}

\chapterTitle{Тема 9. Системи рівнянь}
\noindent{\small\color{gray!80} Розділ 2. Рівняння і нерівності \quad\textbullet\quad Усього завдань: 40 (НМТ \mbox{2023~--- 13}, \mbox{2024~--- 10}, \mbox{2025~--- 6}, \mbox{2026~--- 11}); \mbox{тестових~--- 40}.}\par\vspace{0.2cm}

\typeTitle{Завдання з вибором однієї відповіді (А--Д)}

\begin{samepage}
% Завдання 1
\zadtask{Розв'яжіть систему рівнянь $\begin{cases} 10x - 4y = 26, \\ 6x + 4y = 6. \end{cases}$ Для одержаного розв'язку $(x_0; y_0)$ обчисліть добуток $x_0 \cdot y_0$. \nmtyear{2023}}
\answerTable{$3$}{$-3$}{$6$}{$4$}{$-6$}
\par\penalty-20
\end{samepage}

\begin{samepage}
% Завдання 2
\zadtask{Розв'яжіть систему рівнянь $\begin{cases} 2^{3x-y} = 2^7, \\ 2x + y = 3. \end{cases}$ Якщо $(x_0; y_0)$ --- розв'язок цієї системи, то $x_0 \cdot y_0 =$ \nmtyear{2023}}
\answerTable{$-1$}{$2$}{$-3$}{$1$}{$-2$}
\par\penalty-20
\end{samepage}

\begin{samepage}
% Завдання 3
\zadtask{Розв'яжіть систему рівнянь $\begin{cases} x(4 - y) = 14, \\ x(4 - y) = 3x - 7. \end{cases}$ Якщо $(x_0; y_0)$ --- розв'язок цієї системи, то $y_0 =$ \nmtyear{2023}}
\answerTable{$-6$}{$2$}{$7$}{$6$}{$-2$}
\par\penalty-20
\end{samepage}

\begin{samepage}
% Завдання 4
\zadtask{Розв'яжіть систему рівнянь $\begin{cases} \dfrac{x}{4} = 3y - 1, \\ x - 2y = 1. \end{cases}$ Для одержаного розв'язку $(x_0; y_0)$ системи знайдіть суму $x_0 + y_0$. \nmtyear{2023}}
\answerTable{$2{,}5$}{$1$}{$1{,}6$}{$-2$}{$22$}
\par\penalty-20
\end{samepage}

\begin{samepage}
% Завдання 5
\zadtask{Розв'яжіть систему рівнянь $\begin{cases} x + 2y = 8, \\ x - 2y = -4. \end{cases}$ Якщо $(x_0; y_0)$ --- розв'язок цієї системи, то $x_0 \cdot y_0 =$ \nmtyear{2023}}
\answerTable{$2$}{$3$}{$5$}{$6$}{$-10$}
\par\penalty-20
\end{samepage}

\begin{samepage}
% Завдання 6
\zadtask{Розв'яжіть систему рівнянь $\begin{cases} \dfrac{2x + y}{x} = 5, \\ 4x - y = -2. \end{cases}$ Якщо $(x_0; y_0)$ --- розв'язок цієї системи, то $x_0 + y_0 =$ \nmtyear{2023}}
\answerTable{$-8$}{$4$}{$8$}{$-6$}{$-4$}
\par\penalty-20
\end{samepage}

\begin{samepage}
% Завдання 7
\zadtask{Розв'яжіть систему рівнянь $\begin{cases} \dfrac{3}{x + y} = -\dfrac{1}{2}, \\ 2x - y = 12. \end{cases}$ Якщо $(x_0; y_0)$ --- розв'язок цієї системи, то $x_0 \cdot y_0 =$ \nmtyear{2023}}
\answerTable{$-8$}{$2$}{$-10$}{$16$}{$-16$}
\par\penalty-20
\end{samepage}

\begin{samepage}
% Завдання 8
\zadtask{Розв'яжіть систему рівнянь $\begin{cases} x(y + 2) = 12, \\ x(2y - 1) = 12. \end{cases}$ Якщо $(x_0; y_0)$ --- розв'язок цієї системи, то $x_0 =$ \nmtyear{2023}}
\answerTableTall{$12$}{$2{,}4$}{$4$}{$3$}{$\dfrac{5}{12}$}
\par\penalty-20
\end{samepage}

\begin{samepage}
% Завдання 9
\zadtask{Розв'яжіть систему рівнянь $\begin{cases} \dfrac{1}{3y} = \dfrac{2}{x}, \\ x - y = 30. \end{cases}$ Якщо $(x_0; y_0)$ --- розв'язок цієї системи, то $x_0 + y_0 =$ \nmtyear{2023}}
\answerTable{$42$}{$35$}{$150$}{$36$}{$-150$}
\par\penalty-20
\end{samepage}

\begin{samepage}
% Завдання 10
\zadtask{Розв'яжіть систему рівнянь $\begin{cases} \dfrac{x}{2y} = -3, \\ x + y = 10. \end{cases}$ Якщо $(x_0; y_0)$ --- розв'язок цієї системи, то $x_0 =$ \nmtyear{2023}}
\answerTable{$4$}{$2$}{$12$}{$-2$}{$-12$}
\par\penalty-20
\end{samepage}

\begin{samepage}
% Завдання 11
\zadtask{Розв'яжіть систему рівнянь $\begin{cases} xy - 2x = 3, \\ xy = 2x - y + 11. \end{cases}$ Якщо $(x_0; y_0)$ --- розв'язок цієї системи, то $x_0 \cdot y_0 =$ \nmtyear{2023}}
\answerTable{$8$}{$4$}{$0{,}5$}{$8{,}5$}{$16$}
\par\penalty-20
\end{samepage}

\begin{samepage}
% Завдання 12
\zadtask{Розв'яжіть систему рівнянь $\begin{cases} \sqrt{2x} = \sqrt{6}, \\ x - 4y = 7. \end{cases}$ Якщо $(x_0; y_0)$ --- розв'язок цієї системи, то $y_0 =$ \nmtyear{2023}}
\answerTable{$3$}{$1{,}25$}{$-1$}{$-6$}{$-2{,}5$}
\par\penalty-20
\end{samepage}

\begin{samepage}
% Завдання 13
\zadtask{Розв'яжіть систему рівнянь $\begin{cases} 2x = \dfrac{y}{5}, \\ 4x - y = 3. \end{cases}$ Якщо $(x_0; y_0)$ --- розв'язок цієї системи, то $x_0 + y_0 =$ \nmtyear{2023}}
\answerTable{$2{,}5$}{$-0{,}5$}{$-5{,}5$}{$-5$}{$-4{,}5$}
\par\penalty-20
\end{samepage}

\begin{samepage}
% Завдання 14
\zadtask{Розв'яжіть систему рівнянь $\begin{cases} \dfrac{8}{x} = -2, \\ x + \dfrac{y}{2} = 0. \end{cases}$ \nmtyear{2024}}
\answerTable{$(-0{,}25; 0{,}5)$}{$(4; -8)$}{$(0{,}25; -0{,}5)$}{$(-4; -8)$}{$(-4; 8)$}
\par\penalty-20
\end{samepage}

\begin{samepage}
% Завдання 15
\zadtask{Розв'яжіть систему рівнянь $\begin{cases} \dfrac{2}{y} = \dfrac{7}{x}, \\ x - 4y = 2. \end{cases}$ Якщо $(x_0; y_0)$ --- розв'язок системи, то $x_0 + y_0 =$ \nmtyear{2024}}
\answerTable{$-18$}{$10$}{$-10$}{$22$}{$18$}
\par\penalty-20
\end{samepage}

\begin{samepage}
% Завдання 16
\zadtask{Розв'яжіть систему рівнянь $\begin{cases} x + 2y = 8, \\ 3x - 4y = -1. \end{cases}$ Якщо $(x_0; y_0)$ --- розв'язок системи, то $x_0 + y_0 =$ \nmtyear{2024}}
\answerTable{$2{,}5$}{$3$}{$5$}{$5{,}5$}{$7{,}6$}
\par\penalty-20
\end{samepage}

\begin{samepage}
% Завдання 17
\zadtask{Розв'яжіть систему рівнянь $\begin{cases} 2x - \dfrac{y}{3} = 7, \\[0.6cm] 4x + \dfrac{2y}{3} = 6. \end{cases}$ Якщо $(x_0; y_0)$ --- розв'язок системи, то $x_0 + y_0 =$ \nmtyear{2024}}
\answerTable{$-1{,}5$}{$-6$}{$-3{,}5$}{$2{,}5$}{$-4{,}5$}
\par\penalty-20
\end{samepage}

\begin{samepage}
% Завдання 18
\zadtask{Розв'яжіть систему рівнянь $\begin{cases} \dfrac{2}{x} = -0{,}5, \\[0.6cm] \dfrac{y}{3} = 2. \end{cases}$ \nmtyear{2024}}
\answerTableTall{$\left(-\dfrac{1}{4}; \dfrac{2}{3}\right)$}{$\left(-\dfrac{1}{4}; 6\right)$}{$(6; -4)$}{$\left(-4; \dfrac{2}{3}\right)$}{$(-4; 6)$}
\par\penalty-20
\end{samepage}

\begin{samepage}
% Завдання 19
\zadtask{Розв'яжіть систему рівнянь $\begin{cases} 4x - 3y = -1, \\ 5x - 2y = 4. \end{cases}$ \nmtyear{2024}}
\answerTable{$(-2; 3)$}{$(2; -3)$}{$(3; 2)$}{$(-2; -3)$}{$(2; 3)$}
\par\penalty-20
\end{samepage}

\begin{samepage}
% Завдання 20
\zadtask{Розв'яжіть систему рівнянь $\begin{cases} 2^x + 2y = 0, \\ 2^x - 4y = 6. \end{cases}$ \nmtyear{2024}}
\answerTable{$(-3; 3)$}{$(3; -3)$}{$(1; -1)$}{$(-1; 1)$}{$(2; -2)$}
\par\penalty-20
\end{samepage}

\begin{samepage}
% Завдання 21
\zadtask{Розв'яжіть систему рівнянь $\begin{cases} \dfrac{2}{x} + \dfrac{5}{y} = 5, \\[0.6cm] \dfrac{1}{x} - \dfrac{2}{y} = 7. \end{cases}$ Якщо $(x_0; y_0)$ --- розв'язок системи, то $\dfrac{1}{x_0} + \dfrac{1}{y_0} =$ \nmtyear{2024}}
\answerTable{$4$}{$3$}{$5$}{$6$}{$12$}
\par\penalty-20
\end{samepage}

\begin{samepage}
% Завдання 22
\zadtask{Розв'яжіть систему рівнянь $\begin{cases} y - \dfrac{2}{x} = 3, \\[0.6cm] y + \dfrac{1}{x} = 9. \end{cases}$ Якщо $(x_0; y_0)$ --- розв'язок системи, то $x_0 + y_0 =$ \nmtyear{2024}}
\answerTable{$2$}{$3{,}5$}{$4{,}5$}{$6$}{$7{,}5$}
\par\penalty-20
\end{samepage}

\begin{samepage}
% Завдання 23
\zadtask{Розв'яжіть систему рівнянь $\begin{cases} \sqrt{y} - \dfrac{6}{x} = 6, \\[0.6cm] \sqrt{y} + \dfrac{4}{x} = 1. \end{cases}$ Якщо $(x_0; y_0)$ --- розв'язок системи, то $x_0 + y_0 =$ \nmtyear{2024}}
\answerTable{$2$}{$9$}{$-2$}{$7$}{$83$}
\par\penalty-20
\end{samepage}

\begin{samepage}
% Завдання 24
\zadtask{Розв'яжіть систему рівнянь $\begin{cases} \dfrac{x}{2x + y} = \dfrac{1}{5}, \\[0.6cm] 4x + 2 = y. \end{cases}$ Якщо $(x_0; y_0)$ --- розв'язок системи, то $x_0 + y_0 =$ \nmtyear{2025}}
\answerTable{$12$}{$8$}{$-6$}{$-8$}{$6$}
\par\penalty-20
\end{samepage}

\begin{samepage}
% Завдання 25
\zadtask{Розв'яжіть систему рівнянь $\begin{cases} 2^{y+2x} = 64, \\ 2x - y = 12. \end{cases}$ Якщо $(x_0; y_0)$ --- розв'язок системи, то $x_0 + y_0 =$ \nmtyear{2025}}
\answerTable{$-3$}{$4{,}5$}{$1{,}5$}{$9$}{$7{,}5$}
\par\penalty-20
\end{samepage}

\begin{samepage}
% Завдання 26
\zadtask{Розв'яжіть систему рівнянь $\begin{cases} -2x = 6, \\[0.6cm] \dfrac{y}{3} = 9. \end{cases}$ \nmtyear{2025}}
\answerTable{$(-3; 27)$}{$(3; -27)$}{$(4; 3)$}{$(4; 27)$}{$(-3; 3)$}
\par\penalty-20
\end{samepage}

\begin{samepage}
% Завдання 27
\zadtask{Розв'яжіть систему рівнянь $\begin{cases} x - \dfrac{3}{y} = 3, \\[0.6cm] x + \dfrac{1}{y} = 11. \end{cases}$ Якщо $(x_0; y_0)$ --- розв'язок системи, то $x_0 + y_0 =$ \nmtyear{2025}}
\answerTable{$-2{,}5$}{$2$}{$9{,}5$}{$6{,}5$}{$5$}
\par\penalty-20
\end{samepage}

\begin{samepage}
% Завдання 28
\zadtask{Розв'яжіть систему рівнянь $\begin{cases} \dfrac{y}{4} = 2, \\[0.6cm] x = -2y + 1. \end{cases}$ \nmtyear{2025}}
\answerTable{$(-3; 2)$}{$(-5; 2)$}{$(-15; 8)$}{$(-17; 8)$}{$(0; 0{,}5)$}
\par\penalty-20
\end{samepage}

\begin{samepage}
% Завдання 29
\zadtask{Розв'яжіть систему рівнянь $\begin{cases} 3x + 9 = x + 1, \\ |x - 1| + 2y + 7 = 0. \end{cases}$ Якщо $(x_0; y_0)$ --- розв'язок системи, то $x_0 + y_0 =$ \nmtyear{2025}}
\answerTable{$-10$}{$10{,}5$}{$-0{,}5$}{$-2$}{$-5$}
\par\penalty-20
\end{samepage}

\begin{samepage}
% НМТ 2026, сесія 23.05, завдання №15 --- також у темі: Показникові рівняння
\zadtask{Розв'яжіть систему рівнянь $\begin{cases} 5^x = \sqrt{5}, \\ \dfrac{y + 1}{2x + 3} = \dfrac{1 - 4x}{2}. \end{cases}$ Якщо $(x_0;\, y_0)$~--- розв'язок системи, то $x_0 + y_0 = $ \nmtyear{2026}}
\answerTable{$-2{,}5$}{$-1{,}5$}{$2{,}5$}{$-3{,}5$}{$1{,}5$}
% Відповідь: А
\ifshowsolutions\solution{З першого рівняння $5^x=5^{1/2}$, тож $x_0=0{,}5$. Підставимо: $\dfrac{y+1}{2\cdot 0{,}5+3}=\dfrac{1-4\cdot 0{,}5}{2}$, тобто $\dfrac{y+1}{4}=-\dfrac{1}{2}$, звідки $y+1=-2$, $y_0=-3$. Тоді $x_0+y_0=0{,}5-3=-2{,}5$. Відповідь: \textbf{А}.}\fi
\par\penalty-20
\end{samepage}

\begin{samepage}
% НМТ 2026, сесія 30.05, завдання №1
\zadtask{Розв'яжіть систему рівнянь $\begin{cases} \dfrac{2}{x} = -0{,}5 \\[0.2cm] \dfrac{y}{3} = 2 \end{cases}$ \nmtyear{2026}}
\answerTableTall{$(4;\, 6)$}{$(-4;\, 6)$}{$(-4;\, -6)$}{$(4;\, -6)$}{$\left(-\dfrac{1}{4};\, 6\right)$}
% Відповідь: Б
\ifshowsolutions\solution{З першого рівняння $\dfrac{2}{x}=-0{,}5\Rightarrow x=-4$; з другого $\dfrac{y}{3}=2\Rightarrow y=6$. Розв'язок $(-4;\,6)$. Відповідь: \textbf{Б}.}\fi
\par\penalty-20
\end{samepage}

\begin{samepage}
% НМТ 2026, сесія 2.06, завдання №13
\zadtask{Розв'яжіть систему рівнянь $\begin{cases} xy - 2x = 3, \\ xy - 2x + y = 11. \end{cases}$ \nmtyear{2026}}
\answerTableTall{$\left(\dfrac{1}{2};\, 8\right)$}{$(2;\, 8)$}{$\left(\dfrac{1}{2};\, -8\right)$}{$\left(-\dfrac{1}{2};\, 8\right)$}{$\left(8;\, \dfrac{1}{2}\right)$}
% Відповідь: А
\ifshowsolutions\solution{Віднявши перше рівняння від другого, дістаємо $y=11-3=8$. Підставляємо у перше: $8x-2x=3$, $6x=3$, $x=\dfrac{1}{2}$. Отже $\left(\dfrac{1}{2};\,8\right)$. Відповідь: \textbf{А}.}\fi
\par\penalty-20
\end{samepage}

\begin{samepage}
% НМТ 2026, сесія 3.06, завдання №15
\zadtask{Розв'яжіть систему рівнянь $\begin{cases} \log_2 x = 2, \\[0.15cm] \dfrac{y+10}{x+1} = \dfrac{x+2}{2}. \end{cases}$ Якщо $(x_0;\,y_0)$~--- розв'язок системи, то $x_0 + y_0 = {}$ \nmtyear{2026}}
\answerTable{$18$}{$29$}{$-6$}{$-4$}{$9$}
% Відповідь: Д
\ifshowsolutions\solution{Із $\log_2 x = 2$ маємо $x_0 = 4$. Тоді $\dfrac{y+10}{5} = \dfrac{6}{2} = 3$, звідки $y+10 = 15$, $y_0 = 5$. Отже $x_0 + y_0 = 9$. Відповідь: \textbf{Д}.}\fi
\par\penalty-20
\end{samepage}

\begin{samepage}
% НМТ 2026, сесія 4.06, завдання №4
\zadtask{Розв'яжіть систему рівнянь $\begin{cases} x + 2y = 8, \\ x - 2y = -4. \end{cases}$ \nmtyear{2026}}
\answerTable{$(6;\, 1)$}{$(3;\, 2)$}{$(-2;\, 5)$}{$(2;\, 3)$}{$(2;\, 4)$}
% Відповідь: Г
\ifshowsolutions\solution{Додаємо рівняння: $2x=4$, $x=2$. Тоді $2+2y=8$, $2y=6$, $y=3$. Отже $(2;\,3)$. Відповідь: \textbf{Г}.}\fi
\par\penalty-20
\end{samepage}

\begin{samepage}
% НМТ 2026, сесія 5.06, завдання №4
\zadtask{Розв'яжіть систему рівнянь $\begin{cases} \dfrac{8}{x} = -2, \\[0.2cm] x + \dfrac{y}{2} = 0. \end{cases}$ \nmtyear{2026}}
\answerTable{$(-4;\, 8)$}{$(4;\, -8)$}{$(-4;\, -8)$}{$(4;\, 8)$}{$(-2;\, 4)$}
% Відповідь: А
\ifshowsolutions\solution{З першого рівняння $\dfrac{8}{x}=-2\Rightarrow x=-4$. Підставляємо в друге: $-4+\dfrac{y}{2}=0\Rightarrow y=8$. Розв'язок $(-4;\,8)$. Відповідь: \textbf{А}.}\fi
\par\penalty-20
\end{samepage}

\begin{samepage}
% НМТ 2026, сесія 9.06, завдання №15
\zadtask{Розв'яжіть систему рівнянь
\[
\begin{cases} x + y = -4, \\ xy + y^2 = 12. \end{cases}
\]
Якщо $(x_0; y_0)$~--- розв'язок системи, то $x_0 \cdot y_0 = $ \nmtyear{2026}}
\answerTable{$-1$}{$3$}{$2$}{$-4$}{$-3$}
\par\penalty-20
\end{samepage}

\begin{samepage}
% НМТ 2026, сесія 10.06, завдання №15 --- також у темі: Логарифмічні рівняння
\zadtask{Розв'яжіть систему рівнянь
\[
\begin{cases} \log_2 x + \log_2 y = 2, \\[2pt] x + \dfrac{1}{y} = 10. \end{cases}
\]
Якщо $(x_0; y_0)$~--- розв'язок системи, то $x_0 + 2y_0 =$ \nmtyear{2026}}
\answerTable{$2$}{$8{,}5$}{$9$}{$10$}{$8$}
\par\penalty-20
\end{samepage}

\begin{samepage}
% НМТ 2026, сесія 16.06, завдання №15
\zadtask{Розв'яжіть систему рівнянь
$\begin{cases}
2x + y^3 = -15, \\
4x - y^3 = 51.
\end{cases}$ \\
Якщо $(x_0; y_0)$ --- розв'язок системи, то $y_0 =$ \nmtyear{2026}}
\answerTable{$-3$}{$-27$}{$6$}{$-9$}{$3$}
% Відповідь: А
\par\penalty-20
\end{samepage}

\begin{samepage}
% НМТ 2026, сесія 17.06, завдання №4
\zadtask{Укажіть пару чисел, що є розв'язком системи рівнянь $\begin{cases}xy = 32,\\ y = \dfrac{x}{2}.\end{cases}$ \nmtyear{2026}}
\answerTableTall{$(4; 8)$}{$(16; 2)$}{$(-4; -8)$}{$(2; 4)$}{$(8; 4)$}
% Відповідь: Д
\par\penalty-20
\end{samepage}

\begin{samepage}
% НМТ 2026, сесія 18.06, завдання №15 --- також у темі: Показникові рівняння
\zadtask{Розв'яжіть систему рівнянь $\begin{cases}2^{3x - y} = 2^7,\\ 2x + y = 3.\end{cases}$ Якщо $(x_0;\, y_0)$~--- розв'язок цієї системи, то $x_0 \cdot y_0 = $ \nmtyear{2026}}
\answerTable{$2$}{$-2$}{$-1$}{$-3$}{$1$}
% Відповідь: Б
\par\penalty-20
\end{samepage}

\ifshowsolutions
\sectionTitle{ВІДПОВІДІ ДО ЗАВДАНЬ НМТ 2026}
\begin{multicols}{3}\noindent
\noindent\textbf{30.}~А \ {\scriptsize\color{gray}(23.05, №15)}\par
\noindent\textbf{31.}~Б \ {\scriptsize\color{gray}(30.05, №1)}\par
\noindent\textbf{32.}~А \ {\scriptsize\color{gray}(2.06, №13)}\par
\noindent\textbf{33.}~Д \ {\scriptsize\color{gray}(3.06, №15)}\par
\noindent\textbf{34.}~Г \ {\scriptsize\color{gray}(4.06, №4)}\par
\noindent\textbf{35.}~А \ {\scriptsize\color{gray}(5.06, №4)}\par
\noindent\textbf{38.}~А \ {\scriptsize\color{gray}(16.06, №15)}\par
\noindent\textbf{39.}~Д \ {\scriptsize\color{gray}(17.06, №4)}\par
\noindent\textbf{40.}~Б \ {\scriptsize\color{gray}(18.06, №15)}\par
\end{multicols}
\fi

\endgroup
\end{document}
